% !TEX root = ../diary-book.tex

\section{令和7年09月19日 誤差の見積もりに挑戦}

令和6年10月19日の日記に従って三平面の交点を用いる場合の誤差評価を行う。ベクトルでの微分が使えると楽なので、まずはその準備をする。

行列を使えば交点の式は簡単に解ける。というか形を変えて簡単にかける。

\begin{align*}
    \begin{pmatrix}
        \bm{n}_1^T \\
        \bm{n}_2^T \\
        \bm{n}_3^T
    \end{pmatrix}
    \bm{x}
    &=
    \begin{pmatrix}
        d_1 \\
        d_2 \\
        d_3
    \end{pmatrix} \\
    A\bm{x} &= \bm{d} \\
    \bm{x} &= A^{-1}\bm{d}
\end{align*}

これを下に一次の線形応答を求めると以下のようになる。

\begin{align*}
    \Delta \bm{x} &= \frac{\partial \bm{x}}{\partial A} \Delta A + \frac{\partial \bm{x}}{\partial \bm{d}} \Delta \bm{d} \\
    &= -A^{-1}(\Delta A)A^{-1}\bm{d} + A^{-1}\Delta \bm{d}
\end{align*}

この式の導出の仕方はそれほど単純ではなく、クロネッカー積を使うと楽にできるらしい。

統計による誤算の伝搬は二次の項であり、共分散行列とヤコビアンによって表現される